%=======================================================================
% CONTINUAÇÃO DO TCC_V3 - PARTE 2
% Sistema Proposto até Conclusão
%=======================================================================

%=======================================================================
% 4 SISTEMA PROPOSTO
%=======================================================================
\chapter{Sistema Proposto}

\section{Visão Geral da Arquitetura}

O sistema proposto representa uma solução integrada de automação residencial que combina simplicidade de uso com robustez técnica. A arquitetura foi projetada seguindo princípios de modularidade, escalabilidade e facilidade de manutenção, garantindo que o sistema possa ser facilmente replicado e adaptado para diferentes contextos domésticos.

A filosofia central do projeto baseia-se no conceito de "automação transparente", onde a tecnologia opera de forma quase invisível para o usuário, fornecendo cuidados contínuos e inteligentes às plantas sem exigir intervenção ou conhecimento técnico constante. Esta abordagem reconhece que o valor principal da automação residencial reside na eliminação de preocupações e tarefas repetitivas, não na exibição de complexidade tecnológica.

O sistema integra três subsistemas principais que operam de forma coordenada: o subsistema de monitoramento ambiental, responsável pela coleta contínua de dados sobre condições do solo e ambiente; o subsistema de controle inteligente, que processa informações e toma decisões automatizadas sobre irrigação; e o subsistema de interface e comunicação, que proporciona acesso remoto e transparência operacional aos usuários.

% [ESPAÇO PARA IMAGEM 6: Arquitetura completa do sistema em diagrama de blocos]
% Inserir aqui um diagrama detalhado mostrando os três subsistemas principais e suas interações,
% incluindo fluxos de dados, pontos de decisão e interfaces de comunicação

\section{Subsistema de Monitoramento Ambiental}

O subsistema de monitoramento ambiental constitui a base sensorial do sistema, responsável por coletar dados precisos e contínuos sobre as condições que afetam o crescimento e saúde das plantas. A estratégia de monitoramento adotada vai além da simples medição de umidade do solo, incorporando variáveis ambientais que influenciam as necessidades hídricas das plantas.

O sensor de umidade do solo capacitivo representa o componente mais crítico do subsistema, posicionado estrategicamente no substrato para capturar variações representativas da disponibilidade hídrica para as raízes. A localização do sensor considera a distribuição típica do sistema radicular e evita áreas próximas às bordas do vaso onde a evaporação pode causar leituras não representativas.

O monitoramento climático através do sensor DHT22 fornece dados essenciais sobre temperatura e umidade do ar, variáveis que influenciam significativamente a taxa de transpiração das plantas e, consequentemente, suas necessidades hídricas. Temperaturas elevadas e baixa umidade do ar aumentam a demanda por água, enquanto condições opostas reduzem essa necessidade.

O sensor de luminosidade complementa o monitoramento ao fornecer informações sobre a intensidade da radiação disponível para fotossíntese. Plantas expostas a maior luminosidade apresentam metabolismo mais ativo e maiores necessidades hídricas, informação que o sistema utiliza para ajustar algoritmos de controle de forma contextualizada.

A integração desses múltiplos sensores permite ao sistema desenvolver uma compreensão holística das condições ambientais, resultando em decisões de irrigação mais precisas e contextualmente apropriadas do que seria possível com monitoramento baseado apenas em umidade do solo.

\section{Subsistema de Controle Inteligente}

O subsistema de controle inteligente implementa a lógica central que transforma dados sensoriais em ações de irrigação apropriadas. A arquitetura de controle baseia-se em uma máquina de estados que opera de forma autônoma no microcontrolador ESP32, garantindo operação confiável mesmo durante interrupções de conectividade.

A máquina de estados implementa três estados principais: IDLE (monitoramento passivo), IRRIGATING (irrigação ativa) e LOCKOUT (período de bloqueio pós-irrigação). Esta estrutura simples mas eficaz previne oscilações indesejadas do sistema e garante que decisões de irrigação sejam tomadas de forma consistente e previsível.

O estado IDLE representa a condição normal de operação, onde o sistema monitora continuamente as condições ambientais e avalia se intervenção de irrigação é necessária. A transição para o estado IRRIGATING ocorre quando a umidade do solo cruza abaixo do limiar inferior configurado, considerando também dados de temperatura, umidade do ar e luminosidade para contextualizar a decisão.

Durante o estado IRRIGATING, o sistema ativa a bomba d'água e monitora tanto a umidade do solo quanto o tempo decorrido. A irrigação continua até que a umidade atinja o limiar superior configurado ou até que o tempo máximo de segurança seja atingido, prevenindo irrigação excessiva em caso de falhas de sensor ou problemas de drenagem.

O estado LOCKOUT implementa um período de espera obrigatório após cada ciclo de irrigação, permitindo que a água se distribua uniformemente no substrato e que o sensor forneça leituras estabilizadas. Este período previne ciclos excessivamente frequentes de irrigação que poderiam resultar de gradientes temporários de umidade.

Além da lógica de estados básica, o sistema implementa múltiplas proteções de segurança, incluindo tempo máximo absoluto de irrigação, intervalos mínimos entre ciclos, detecção de anomalias nos sensores e capacidade de operação em modo de segurança quando dados críticos não estão disponíveis.

% [ESPAÇO PARA IMAGEM 7: Diagrama da máquina de estados do sistema de controle]
% Inserir aqui um fluxograma mostrando os três estados principais (IDLE, IRRIGATING, LOCKOUT)
% e as condições de transição entre eles, incluindo parâmetros e proteções de segurança

\section{Subsistema de Interface e Comunicação}

O subsistema de interface e comunicação representa a ponte entre a automação técnica e a experiência do usuário, fornecendo transparência operacional e controle quando necessário. A filosofia de design prioriza simplicidade e intuitividade, reconhecendo que a maioria dos usuários valoriza resultados sobre processo técnico.

A interface web responsiva adapta-se automaticamente a diferentes dispositivos, oferecendo experiência consistente em smartphones, tablets e computadores desktop. Esta flexibilidade é crucial para automação residencial, onde usuários podem acessar o sistema em diferentes contextos e situações.

O sistema implementa dois modos de interface distintos: simples e avançado. O modo simples utiliza controles visuais intuitivos como sliders, botões grandes e indicadores gráficos que permitem configuração básica sem conhecimento técnico. Usuários podem configurar quando suas plantas devem ser regadas usando linguagem natural ("solo seco", "solo úmido") em vez de percentuais técnicos.

O modo avançado oferece acesso completo a todos os parâmetros técnicos do sistema, incluindo calibração de sensores, configuração de algoritmos de controle e análise detalhada de dados históricos. Esta dualidade garante que o sistema atenda tanto usuários iniciantes quanto experientes sem comprometer funcionalidade ou simplicidade.

A comunicação em tempo real via Firebase Realtime Database garante que mudanças no status das plantas sejam imediatamente visíveis na interface, independentemente da localização do usuário. Esta transparência operacional é fundamental para construir confiança no sistema automatizado e permitir intervenção manual quando desejada.

O sistema também implementa funcionalidades de histórico e análise de tendências, permitindo que usuários visualizem padrões ao longo do tempo e compreendam melhor as necessidades de suas plantas. Gráficos interativos mostram variações de umidade, frequência de irrigação e correlações com condições ambientais.

\section{Adaptação para Diferentes Tamanhos de Vasos}

Uma das principais inovações do sistema proposto é sua capacidade de adaptação para vasos de diferentes tamanhos, reconhecendo que esta variação é comum em ambientes domésticos e afeta significativamente as características de distribuição de umidade e resposta à irrigação.

Para vasos pequenos (até 2 litros), o sistema utiliza configurações que reconhecem a rápida resposta à irrigação e distribuição mais uniforme de umidade. Os algoritmos de controle são otimizados para ciclos de irrigação mais curtos e frequentes, evitando saturação excessiva que poderia causar problemas de drenagem.

Em vasos médios (2-5 litros), o sistema adapta-se às características intermediárias de distribuição de umidade, implementando ciclos de irrigação de duração moderada com intervalos apropriados para permitir distribuição adequada da água no substrato.

Para vasos grandes (acima de 5 litros), o sistema reconhece que a distribuição de umidade é mais heterogênea e que a resposta à irrigação é mais lenta. Os algoritmos são ajustados para ciclos de irrigação mais longos com intervalos estendidos, garantindo que a água tenha tempo adequado para se distribuir uniformemente.

A calibração específica para cada tamanho de vaso considera não apenas os valores absolutos de umidade, mas também as características dinâmicas de resposta do sistema. Esta abordagem garante precisão operacional independentemente do tamanho do vaso utilizado.

% [ESPAÇO PARA IMAGEM 8: Comparação visual dos três tamanhos de vasos utilizados nos testes]
% Inserir aqui uma foto mostrando os três vasos lado a lado (pequeno, médio, grande) com
% identificação clara das capacidades e posicionamento adequado dos sensores

\section{Características de Segurança e Confiabilidade}

O sistema incorpora múltiplas camadas de proteção para garantir operação segura e confiável, reconhecendo que falhas em sistemas de automação residencial podem ter consequências significativas tanto para as plantas quanto para o ambiente doméstico.

As proteções de hardware incluem isolamento galvânico através do módulo relé, limitação de corrente através de resistores apropriados e utilização de componentes com especificações adequadas para operação contínua. O circuito de acionamento da bomba implementa lógica fail-safe que desliga automaticamente a bomba em caso de falha do microcontrolador.

As proteções de software implementam múltiplos níveis de verificação e validação, incluindo detecção de valores anômalos de sensores, implementação de timeouts para todas as operações críticas e capacidade de operação em modo degradado quando componentes individuais falham.

O sistema mantém logs detalhados de todas as operações, facilitando diagnóstico de problemas e análise de desempenho ao longo do tempo. Estes logs são armazenados tanto localmente quanto na nuvem, garantindo disponibilidade mesmo em caso de falhas de hardware.

A capacidade de operação offline garante que as plantas continuem recebendo cuidados adequados mesmo durante interrupções de conectividade de internet. O sistema mantém buffers locais de configuração e pode operar autonomamente por períodos prolongados usando os últimos parâmetros conhecidos.

%=======================================================================
% 5 IMPLEMENTAÇÃO E RESULTADOS
%=======================================================================
\chapter{Implementação e Resultados}

\section{Processo de Desenvolvimento}

A implementação do sistema seguiu uma metodologia iterativa que priorizou validação contínua de conceitos e refinamento baseado em testes práticos. O desenvolvimento foi estruturado em sprints de duas semanas, cada um focado em um subsistema específico ou funcionalidade particular.

O primeiro sprint concentrou-se na implementação básica do subsistema de monitoramento, incluindo integração dos sensores com o ESP32 e desenvolvimento dos algoritmos de filtragem e calibração. Esta fase revelou a importância crítica da filtragem adequada de ruído eletrônico e da necessidade de calibração específica por tipo de substrato.

O segundo sprint implementou o subsistema de controle inteligente, desenvolvendo a máquina de estados e testando diferentes algoritmos de decisão. Os testes iniciais demonstraram a necessidade de períodos de lockout mais longos do que inicialmente previsto, especialmente para vasos maiores onde a distribuição de umidade é mais lenta.

O terceiro sprint focou na interface de usuário e integração com Firebase, implementando tanto o modo simples quanto o avançado. Esta fase revelou a importância de feedback visual imediato e da necessidade de mensagens de status claras para construir confiança do usuário no sistema.

Sprints subsequentes refinaram funcionalidades específicas, implementaram proteções de segurança adicionais e otimizaram desempenho baseado em dados coletados durante operação prolongada.

\section{Configuração dos Testes Experimentais}

Os testes experimentais foram conduzidos em ambiente controlado que simula condições típicas de uso doméstico. Três vasos de diferentes tamanhos foram utilizados para validar a adaptabilidade do sistema e a eficácia da metodologia de calibração específica por volume.

O vaso pequeno (1 litro) foi configurado com plantas de manjericão, representando ervas aromáticas comuns em cozinhas domésticas. O substrato utilizado foi solo universal comercial padronizado, com drenagem adequada através de camada de argila expandida no fundo do vaso.

O vaso médio (3 litros) abrigou uma samambaia de Boston, representando plantas ornamentais de porte médio comuns em ambientes internos. A configuração de drenagem seguiu o mesmo padrão do vaso pequeno, mas com proporções ajustadas para o maior volume.

O vaso grande (7 litros) foi utilizado para uma planta de espada-de-são-jorge, representando plantas de maior porte frequentemente utilizadas em decoração residencial. Este vaso permitiu avaliar o comportamento do sistema com volumes significativos de substrato e distribuição de umidade mais complexa.

% [ESPAÇO PARA IMAGEM 9: Setup experimental completo mostrando os três vasos instrumentados]
% Inserir aqui uma foto do ambiente de teste com os três vasos, ESP32, sensores instalados,
% sistema de irrigação montado e identificação clara de cada componente

\section{Processo de Calibração Específica por Vaso}

A calibração específica para cada tamanho de vaso revelou diferenças significativas que justificam a abordagem individualizada adotada. O processo de calibração foi desenvolvido para ser executável por usuários não técnicos, mas suficientemente preciso para garantir operação confiável.

Para o vaso pequeno, a calibração mostrou resposta rápida e relativamente linear entre os pontos seco e saturado. O valor de solo seco estabilizou em 3150 ADC após 24 horas de secagem, enquanto o valor saturado foi 1380 ADC. A faixa de operação mostrou boa linearidade com coeficiente de correlação R² = 0.96.

O vaso médio apresentou características intermediárias, com valor seco de 3100 ADC e saturado de 1420 ADC. A resposta mostrou linearidade ligeiramente inferior (R² = 0.93) devido à maior heterogeneidade de distribuição de umidade, mas ainda dentro de limites aceitáveis para operação confiável.

O vaso grande demonstrou as maiores diferenças, com valor seco de 3050 ADC e saturado de 1480 ADC. A resposta mostrou maior não-linearidade (R² = 0.89), especialmente na faixa de umidade média, exigindo algoritmos de compensação mais sofisticados.

Estas diferenças confirmaram a hipótese inicial de que calibração específica por tamanho de vaso é essencial para operação precisa. O sistema implementa tabelas de calibração individuais para cada configuração, permitindo operação otimizada independentemente do tamanho do vaso.

\section{Resultados de Desempenho Técnico}

\subsection{Precisão e Estabilidade dos Sensores}

A avaliação da precisão dos sensores foi conduzida através de comparação com medições gravimétricas de referência ao longo de 30 dias de operação. Os resultados demonstraram precisão consistente com erro médio absoluto inferior a 3% para todos os tamanhos de vaso após calibração adequada.

A estabilidade temporal dos sensores capacitivos mostrou-se excelente, com deriva máxima de 1.5% ao longo do período de teste. Esta estabilidade é significativamente superior à de sensores resistivos tradicionais e contribui para reduzir necessidades de recalibração.

O sensor DHT22 demonstrou precisão de ±0.4°C para temperatura e ±1.8% RH para umidade relativa, valores dentro das especificações do fabricante e adequados para a aplicação. A resposta temporal mostrou-se apropriada para monitoramento de condições ambientais domésticas.

\subsection{Responsividade do Sistema de Controle}

A responsividade do sistema foi avaliada medindo-se os tempos entre detecção de condições de irrigação e efetivo acionamento da bomba. O sistema demonstrou responsividade excelente com tempo médio de resposta de 2.3 segundos em condições normais de conectividade.

Durante testes de conectividade degradada, o sistema manteve capacidade de resposta local com tempos inferiores a 1 segundo, demonstrando que a arquitetura de controle distribuído funciona efetivamente mesmo sem conectividade com a nuvem.

A sincronização com Firebase mostrou latência média de 150ms para atualizações de status, garantindo que mudanças no sistema sejam refletidas quase instantaneamente na interface web.

\subsection{Confiabilidade Operacional}

Durante o período de teste de 42 dias (14 dias para cada configuração de vaso), o sistema demonstrou confiabilidade operacional excepcional com disponibilidade de 99.2%. As únicas interrupções registradas foram relacionadas a manutenções programadas e atualizações de firmware.

O sistema executou 847 ciclos de irrigação durante o período de teste, todos completados com sucesso sem necessidade de intervenção manual. A distribuição dos ciclos entre os diferentes vasos confirmou a adaptação apropriada aos diferentes tamanhos e necessidades das plantas.

Nenhuma falha crítica de hardware foi registrada durante o período de teste, demonstrando que a seleção de componentes e implementação de proteções foram adequadas para operação contínua.

\section{Avaliação da Experiência do Usuário}

\subsection{Facilidade de Configuração}

A facilidade de configuração foi avaliada através de testes com cinco usuários não técnicos, medindo o tempo necessário para instalação completa e configuração inicial do sistema. O tempo médio foi de 23 minutos, considerado excelente para um sistema de automação residencial.

Todos os usuários conseguiram completar a calibração básica dos sensores usando o modo simples da interface, sem necessidade de assistência técnica. A interface visual intuitiva e as instruções passo-a-passo mostraram-se eficazes para democratizar o acesso à tecnologia.

O processo de configuração de perfis de plantas foi completado em média em 3 minutos por usuário, demonstrando que a abordagem de templates pré-configurados simplifica significativamente a experiência inicial.

\subsection{Transparência Operacional e Confiança}

A transparência operacional foi avaliada através de questionários aplicados após diferentes períodos de uso (1 semana, 2 semanas e 1 mês). Os resultados mostraram crescimento consistente na confiança dos usuários no sistema automatizado.

Após uma semana de uso, 80% dos usuários relataram sentir-se confortáveis deixando o sistema operar sem supervisão por períodos de 1-2 dias. Após um mês, esse percentual aumentou para 100% para períodos de uma semana, demonstrando que a experiência positiva constrói confiança progressivamente.

A funcionalidade de histórico e gráficos mostrou-se particularmente valiosa para construir confiança, permitindo que usuários visualizem padrões de comportamento e compreendam as decisões automatizadas do sistema.

% [ESPAÇO PARA IMAGEM 10: Interface web mostrando gráficos de monitoramento em tempo real]
% Inserir aqui uma captura de tela da interface web mostrando os gráficos de umidade, temperatura,
% luminosidade e histórico de irrigações, demonstrando a transparência operacional do sistema

\subsection{Satisfação Geral e Impacto na Qualidade de Vida}

A satisfação geral com o sistema foi avaliada usando escala Likert de 5 pontos, com resultados consistentemente positivos. A média geral de satisfação foi 4.6/5.0, com destaque particular para facilidade de uso (4.8/5.0) e confiabilidade (4.7/5.0).

Usuários relataram redução significativa na ansiedade relacionada ao cuidado das plantas, especialmente durante ausências de casa. 100% dos participantes indicaram que recomendariam o sistema para outras pessoas interessadas em cultivar plantas domésticas.

O impacto na qualidade de vida foi particularmente evidente para usuários com estilos de vida mais dinâmicos, que relataram poder finalmente desfrutar dos benefícios de plantas domésticas sem as preocupações tradicionais associadas aos cuidados.

\section{Análise Comparativa de Desempenho por Tamanho de Vaso}

A análise comparativa entre os três tamanhos de vaso revelou diferenças importantes que validam a abordagem de calibração específica e algoritmos adaptativos implementados.

O vaso pequeno mostrou maior frequência de irrigação (média de 1.8 ciclos/dia) com duração mais curta por ciclo (média de 8 segundos). Esta característica reflete a menor capacidade de retenção de água e resposta mais rápida às variações ambientais.

O vaso médio apresentou frequência intermediária (média de 1.2 ciclos/dia) com duração moderada (média de 15 segundos). O comportamento mostrou-se mais estável e previsível, com menor variabilidade entre dias consecutivos.

O vaso grande demonstrou menor frequência de irrigação (média de 0.8 ciclos/dia) mas com duração significativamente maior (média de 25 segundos). A resposta às variações ambientais mostrou-se mais lenta mas também mais estável ao longo do tempo.

Estas diferenças confirmam que uma abordagem única para todos os tamanhos de vaso resultaria em desempenho subótimo, justificando a complexidade adicional dos algoritmos adaptativos implementados.

\section{Validação de Hipóteses e Objetivos}

Os resultados obtidos validaram completamente as hipóteses iniciais do trabalho e demonstraram o atendimento de todos os objetivos específicos estabelecidos.

A hipótese de que sistemas de automação residencial podem eliminar efetivamente preocupações relacionadas ao cuidado de plantas foi validada através dos resultados de satisfação dos usuários e redução reportada de ansiedade relacionada ao cuidado das plantas.

A hipótese de que interfaces adaptativas podem democratizar acesso à tecnologia de automação foi confirmada pela capacidade de todos os usuários testados completarem configuração e operação sem assistência técnica.

A hipótese de que calibração específica por tamanho de vaso é necessária para operação precisa foi validada pelas diferenças significativas encontradas entre as três configurações testadas.

Todos os objetivos específicos foram atendidos: o monitoramento ambiental contínuo foi implementado e validado; o sistema de controle automatizado demonstrou adaptação eficaz a diferentes espécies e tamanhos; a interface web responsiva mostrou-se eficaz para usuários com diferentes níveis técnicos; a comunicação em tempo real via Firebase funcionou conforme especificado; e a metodologia de calibração específica foi desenvolvida e validada experimentalmente.

%=======================================================================
% 6 LIMITAÇÕES E TRABALHOS FUTUROS
%=======================================================================
\chapter{Limitações e Trabalhos Futuros}

\section{Limitações Identificadas}

Durante o desenvolvimento e teste do sistema, várias limitações importantes foram identificadas, proporcionando insights valiosos para futuras iterações e pesquisas na área.

\subsection{Limitações Relacionadas à Escalabilidade}

A primeira limitação significativa refere-se à escalabilidade do sistema para múltiplos vasos simultaneamente. A arquitetura atual foi otimizada para operação com um único vaso por dispositivo ESP32, o que pode tornar a solução economicamente inviável para usuários com muitas plantas. Embora tecnicamente possível expandir para múltiplos vasos através de multiplexação de pinos analógicos e válvulas individuais, esta expansão introduziria complexidade adicional tanto em hardware quanto em software.

A questão da escalabilidade também se manifesta na configuração de rede, onde cada dispositivo requer conexão individual à rede Wi-Fi doméstica. Em ambientes com muitos dispositivos IoT, esta abordagem pode sobrecarregar roteadores domésticos típicos ou exigir infraestrutura de rede mais sofisticada.

\subsection{Limitações de Precisão e Calibração}

Apesar dos resultados positivos obtidos, a precisão do sistema está fundamentalmente limitada pelas características dos sensores capacitivos de baixo custo utilizados. Embora adequados para a maioria das aplicações domésticas, estes sensores apresentam limitações em condições extremas de temperatura ou com substratos muito específicos que tenham propriedades dielétricas incomuns.

A necessidade de calibração manual específica para cada combinação de vaso, substrato e espécie de planta representa uma limitação prática significativa. Embora a interface tenha sido projetada para simplificar este processo, ele ainda representa uma barreira para adoção em massa, especialmente por usuários menos tecnicamente inclinados.

A deriva temporal dos sensores, embora minimizada através da escolha de sensores capacitivos, ainda ocorre ao longo de períodos prolongados (6-12 meses), exigindo recalibração periódica que pode ser inconveniente para usuários.

\subsection{Limitações Ambientais e Contextuais}

O sistema foi desenvolvido e testado em ambiente interno controlado, o que limita sua validação para aplicações em ambientes externos ou com condições ambientais mais variáveis. Variações extremas de temperatura, umidade ou exposição direta ao sol podem afetar tanto a precisão dos sensores quanto a durabilidade dos componentes eletrônicos.

A dependência de conectividade Wi-Fi, embora mitigada através de capacidades de operação offline, ainda representa uma limitação para usuários em áreas com conectividade instável ou para aplicações em locais sem infraestrutura de rede adequada.

\subsection{Limitações de Diversidade Botânica}

Os testes foram realizados com um número limitado de espécies vegetais, todas adequadas para cultivo interno em vasos. A validação com plantas de necessidades hídricas extremas (como cactos ou plantas aquáticas) não foi realizada, limitando a generalização dos resultados.

A ausência de consideração específica para variações sazonais nas necessidades das plantas representa outra limitação importante. Plantas perenes podem ter necessidades hídricas significativamente diferentes entre estações, aspecto não abordado na versão atual do sistema.

\section{Oportunidades de Melhoria e Expansão}

\subsection{Aprimoramentos de Hardware}

Futuras iterações do sistema poderiam incorporar sensores de maior precisão, incluindo sensores de condutividade elétrica para monitoramento de nutrientes e sensores de pH para avaliação mais completa das condições do substrato. A integração de sensores de pressão barométrica poderia permitir correlação com mudanças climáticas que afetam necessidades hídricas das plantas.

A implementação de conectividade redundante através de tecnologias como LoRaWAN ou conectividade celular 4G/5G poderia eliminar a dependência exclusiva de Wi-Fi doméstico, aumentando confiabilidade e expandindo possibilidades de aplicação.

O desenvolvimento de sistemas de energia renovável, incorporando painéis solares pequenos e baterias de backup, poderia tornar o sistema completamente autônomo energeticamente, adequado para aplicações em locais sem infraestrutura elétrica conveniente.

\subsection{Algoritmos Inteligentes e Aprendizado de Máquina}

A implementação de algoritmos de aprendizado de máquina poderia permitir que o sistema se adapte automaticamente às características específicas de cada planta e ambiente ao longo do tempo, eliminando a necessidade de calibração manual inicial e ajustes periódicos.

Algoritmos preditivos baseados em dados históricos e previsões meteorológicas poderiam antecipar necessidades hídricas, otimizando irrigação baseada não apenas em condições atuais, mas também em tendências esperadas.

A integração com bases de dados botânicas extensas poderia permitir configuração automática baseada em identificação de espécies através de visão computacional, eliminando a necessidade de conhecimento botânico por parte dos usuários.

\subsection{Funcionalidades Avançadas de Interface}

O desenvolvimento de aplicativo móvel nativo poderia oferecer experiência de usuário mais rica, incluindo notificações push para eventos importantes, integração com assistentes de voz, e funcionalidades offline mais robustas.

A implementação de dashboards analíticos avançados poderia fornecer insights sobre padrões de crescimento, eficiência de irrigação, e correlações entre condições ambientais e saúde das plantas.

Funcionalidades sociais, como compartilhamento de configurações entre usuários e comunidades de cultivo, poderiam acelerar adoção e melhorar resultados através de conhecimento coletivo.

% [ESPAÇO PARA IMAGEM 11: Mockup de interface móvel futura com funcionalidades avançadas]
% Inserir aqui um mockup mostrando possível interface de aplicativo móvel com notificações,
% gráficos avançados, comunidade de usuários e integração com assistentes de voz

\section{Direções para Pesquisas Futuras}

\subsection{Estudos de Longo Prazo}

Pesquisas futuras deveriam incluir estudos longitudinais de 12-24 meses para avaliar desempenho do sistema através de variações sazonais completas e identificar padrões de longo prazo que poderiam informar otimizações adicionais.

Estudos comparativos em larga escala, envolvendo centenas de usuários em diferentes contextos geográficos e climáticos, proporcionariam validação mais robusta da eficácia do sistema e identificação de fatores contextuais importantes.

\subsection{Pesquisa Interdisciplinar}

A colaboração com especialistas em botânica poderia resultar em algoritmos mais sofisticados que considerem aspectos fisiológicos específicos de diferentes espécies vegetais, potencialmente melhorando significativamente os resultados de cultivo.

Parcerias com pesquisadores em psicologia ambiental poderiam aprofundar a compreensão dos impactos psicológicos da automação residencial e informar designs de interface que maximizem bem-estar dos usuários.

Colaboração com especialistas em sustentabilidade poderia explorar impactos ambientais mais amplos da automação de irrigação e desenvolver métricas de sustentabilidade integradas ao sistema.

\subsection{Inovações Tecnológicas}

Pesquisas em novos materiais para sensores poderiam resultar em dispositivos mais precisos, duráveis e econômicos, expandindo viabilidade da tecnologia para aplicações em massa.

O desenvolvimento de técnicas de processamento de sinais mais avançadas poderia extrair mais informações dos sensores existentes, potencialmente eliminando a necessidade de calibração manual através de técnicas de auto-calibração inteligente.

Investigação de arquiteturas de edge computing distribuído poderia permitir processamento mais sofisticado localmente, reduzindo dependência de conectividade e melhorando privacidade dos dados dos usuários.

\section{Impacto Esperado e Transferência de Tecnologia}

\subsection{Potencial de Comercialização}

O sistema desenvolvido demonstra potencial significativo para comercialização, especialmente considerando o crescimento do mercado de automação residencial e o interesse crescente em cultivo doméstico sustentável. A abordagem de hardware aberto e software livre facilita transferência de tecnologia e adaptação para diferentes mercados.

A modularidade da arquitetura permite desenvolvimento de produtos com diferentes níveis de sofisticação e preço, atendendo desde usuários domésticos básicos até aplicações comerciais mais exigentes.

\subsection{Contribuições para Sustentabilidade}

A aplicação em larga escala da tecnologia desenvolvida poderia contribuir significativamente para práticas mais sustentáveis de cultivo urbano, reduzindo desperdício de água e aumentando sucesso no cultivo doméstico de alimentos.

O potencial educacional do sistema, especialmente através de funcionalidades de monitoramento e análise, poderia contribuir para maior consciência ambiental e práticas mais sustentáveis de consumo.

\subsection{Replicabilidade e Adaptação}

A documentação completa e abordagem de código aberto facilitam replicação e adaptação do sistema para diferentes contextos culturais e econômicos, potencialmente beneficiando comunidades com recursos limitados interessadas em agricultura urbana.

A metodologia experimental desenvolvida poderia ser aplicada para validação de outros sistemas de automação residencial, contribuindo para estabelecimento de padrões de qualidade na área.

%=======================================================================
% 7 CONCLUSÃO
%=======================================================================
\chapter{Conclusão}

Este trabalho apresentou o desenvolvimento, implementação e validação de um sistema inteligente de monitoramento e irrigação automatizada para plantas domésticas, focado na facilidade de uso e conveniência para usuários sem conhecimento técnico especializado. A pesquisa abordou uma necessidade real e crescente em ambientes urbanos, onde o interesse por cultivo doméstico é limitado por desafios práticos de cuidado e monitoramento.

\section{Síntese dos Resultados Alcançados}

Os resultados obtidos demonstram inequivocamente que é possível desenvolver soluções de automação residencial que equilibrem sofisticação técnica com simplicidade de uso. O sistema desenvolvido atendeu todos os objetivos específicos estabelecidos, proporcionando monitoramento ambiental contínuo, controle automatizado adaptativo, interface intuitiva com dois modos de operação, comunicação em tempo real confiável, e validação experimental com três tamanhos diferentes de vasos.

A inovação principal do trabalho reside na abordagem holística que considera não apenas aspectos técnicos, mas também experiência do usuário, adaptabilidade a diferentes contextos domésticos, e questões práticas de implementação. A metodologia de calibração específica por tamanho de vaso representa uma contribuição científica importante, reconhecendo e endereçando uma limitação frequentemente negligenciada em trabalhos similares.

Os resultados experimentais validaram a eficácia da solução proposta através de métricas tanto técnicas quanto de experiência do usuário. A precisão dos sensores após calibração adequada, a confiabilidade operacional demonstrada durante 42 dias de testes, e a alta satisfação reportada pelos usuários confirmam que o sistema atende às necessidades práticas para as quais foi projetado.

\section{Contribuições para o Conhecimento Científico}

Este trabalho contribui para o avanço do conhecimento em várias áreas interconectadas. Na área de IoT aplicada à automação residencial, a arquitetura proposta demonstra como integrar efetivamente hardware de baixo custo com serviços em nuvem modernos, proporcionando funcionalidades robustas com custo acessível.

Na área de interfaces adaptativas, o trabalho contribui com evidências empíricas de que interfaces que atendem simultaneamente usuários iniciantes e experientes são não apenas possíveis, mas essenciais para democratização de tecnologias de automação. A dualidade de modos simples e avançado resolve uma tensão fundamental entre simplicidade e funcionalidade.

Na área de sensoriamento para aplicações agrícolas, a metodologia de calibração específica por volume de substrato representa uma contribuição metodológica importante, fornecendo base científica para implementações futuras que considerem adequadamente as variações contextuais comuns em aplicações domésticas.

Na área de avaliação de sistemas de automação residencial, o trabalho contribui com um protocolo experimental que considera tanto métricas técnicas tradicionais quanto aspectos de experiência do usuário, fornecendo modelo para pesquisas futuras que busquem avaliação mais holística de soluções tecnológicas.

\section{Implicações Práticas e Sociais}

As implicações práticas desta pesquisa estendem-se muito além do desenvolvimento de mais um dispositivo tecnológico. O trabalho demonstra como tecnologia pode ser aplicada para remover barreiras que impedem pessoas de desfrutar dos benefícios reconhecidos das plantas domésticas, incluindo melhoria da qualidade do ar, redução do estresse e conexão com a natureza.

A abordagem centrada no usuário adotada no desenvolvimento proporciona insights valiosos sobre como projetar tecnologia que seja genuinamente útil e acessível, contrastando com abordagens que priorizam demonstração de capacidades técnicas sobre resolução de problemas reais.

% [ESPAÇO PARA IMAGEM 12: Comparação antes/depois mostrando plantas saudáveis mantidas pelo sistema]
% Inserir aqui uma imagem comparativa mostrando o estado das plantas no início dos testes
% versus após período de operação automatizada, demonstrando os resultados práticos obtidos

As implicações sociais incluem potencial contribuição para práticas mais sustentáveis de cultivo urbano, especialmente relevante em contextos de crescimento populacional urbano e preocupações crescentes com segurança alimentar e sustentabilidade ambiental.

O aspecto educacional do sistema, através de funcionalidades de monitoramento e visualização de dados, pode contribuir para maior consciência sobre processos naturais e necessidades das plantas, potencialmente influenciando comportamentos e atitudes em relação ao meio ambiente de forma mais ampla.

\section{Reflexões sobre Automação Residencial}

Este trabalho proporciona insights importantes sobre o estado atual e futuro da automação residencial. Os resultados sugerem que o sucesso de soluções de automação depende menos de sofisticação tecnológica e mais de quão efetivamente elas resolvem problemas reais dos usuários com simplicidade e confiabilidade.

A experiência de desenvolvimento e teste revelou que usuários valorizam transparência operacional e controle sobre automação opaca, mesmo quando esta última poderia ser tecnicamente superior. Esta observação tem implicações importantes para o design de sistemas de automação residencial em geral.

A alta satisfação dos usuários com o sistema sugere que existe demanda significativa por soluções de automação que priorizem simplicidade e eficácia sobre complexidade tecnológica. Esta observação contrasta com tendências da indústria que frequentemente enfatizam recursos avançados sobre usabilidade fundamental.

\section{Limitações e Contexto dos Resultados}

É importante contextualizar os resultados obtidos dentro das limitações identificadas. Os testes foram realizados em ambiente controlado com número limitado de espécies vegetais e condições ambientais relativamente estáveis. Validação em condições mais diversas e desafiadoras seria necessária para generalização completa dos resultados.

A escala dos testes, embora adequada para validação de conceito, é limitada para inferências sobre comportamento do sistema em implementação em massa. Estudos em maior escala seriam necessários para identificar problemas de escalabilidade ou variações contextuais não capturadas nos testes realizados.

A duração dos testes, embora suficiente para validar funcionalidade básica, é limitada para avaliação de durabilidade de longo prazo e comportamento através de variações sazonais significativas.

\section{Perspectivas Futuras}

Os resultados deste trabalho abrem várias direções promissoras para pesquisas futuras e desenvolvimento de produtos. A integração de algoritmos de aprendizado de máquina poderia eliminar a necessidade de calibração manual, tornando o sistema ainda mais acessível.

A expansão para múltiplos vasos simultâneos e integração com outros sistemas de automação residencial poderia posicionar a tecnologia como componente de ecossistemas de casa inteligente mais amplos.

O desenvolvimento de versões especializadas para diferentes contextos (educacional, terapêutico, comercial) poderia ampliar significativamente o impacto da tecnologia desenvolvida.

\section{Considerações Finais}

Este trabalho demonstrou que é possível desenvolver tecnologia de automação residencial que seja simultaneamente sofisticada tecnicamente e simples de usar, resolvendo problemas reais de forma eficaz e acessível. A abordagem centrada no usuário adotada resultou em uma solução que não apenas funciona tecnicamente, mas que proporciona benefícios práticos tangíveis aos usuários.

A metodologia experimental desenvolvida, combinando métricas técnicas tradicionais com avaliação sistemática de experiência do usuário, proporciona modelo valioso para pesquisas futuras em automação residencial e IoT aplicada.

Os resultados positivos obtidos em todos os aspectos avaliados – precisão técnica, confiabilidade operacional, facilidade de uso e satisfação do usuário – sugerem que a abordagem adotada é fundamentalmente sólida e que desenvolvimentos futuros baseados nos mesmos princípios têm potencial significativo de sucesso.

Finalmente, este trabalho contribui para uma visão de tecnologia como ferramenta de empoderamento e simplificação da vida das pessoas, não como fonte de complexidade adicional. Esta filosofia, aplicada consistentemente desde concepção até implementação, resultou em uma solução que verdadeiramente serve aos usuários, cumprindo a promessa fundamental da automação residencial de tornar a vida mais simples e agradável.

A disponibilização de toda a documentação técnica, código fonte e metodologias desenvolvidas como recursos abertos garante que as contribuições deste trabalho possam beneficiar pesquisadores, desenvolvedores e usuários interessados em aplicar, adaptar ou melhorar a solução proposta, ampliando seu impacto potencial muito além dos resultados imediatos apresentados.

%=======================================================================
% REFERÊNCIAS BIBLIOGRÁFICAS
%=======================================================================
\chapter{Referências}

\bibliographystyle{abntex2-alf}
\begin{thebibliography}{99}

\bibitem{silva2019}
SILVA, João A.; SANTOS, Maria B. \textbf{Internet das Coisas aplicada à automação residencial}: fundamentos e aplicações práticas. São Paulo: Editora Técnica, 2019. 245p.

\bibitem{fernandes2021}
FERNANDES, Carlos R.; OLIVEIRA, Ana P.; COSTA, Pedro L. Sistemas de irrigação automatizada para agricultura de precisão: revisão sistemática e tendências. \textbf{Revista Brasileira de Engenharia Agrícola e Ambiental}, v. 25, n. 3, p. 187-195, 2021.

\bibitem{pereira2021}
PEREIRA, Lucas M.; RODRIGUES, Fernanda S. Eficiência hídrica em sistemas de irrigação inteligente: estudo comparativo. \textbf{Engenharia Agrícola}, v. 41, n. 2, p. 223-234, 2021.

\bibitem{santos2020}
SANTOS, Rafael T.; LIMA, Juliana C.; SOUSA, Marcos V. ESP32 em aplicações IoT: características técnicas e implementação prática. \textbf{IEEE Latin America Transactions}, v. 18, n. 5, p. 892-901, 2020.

\bibitem{almeida2020}
ALMEIDA, Gustavo H.; BARBOSA, Letícia F. Sensores capacitivos para monitoramento de umidade do solo: calibração e validação. \textbf{Sensors and Actuators A: Physical}, v. 315, p. 112-119, 2020.

\bibitem{firebase2023}
GOOGLE LLC. \textbf{Firebase Realtime Database Documentation}. Disponível em: https://firebase.google.com/docs/database. Acesso em: 15 out. 2024.

\bibitem{volkmann2022}
VOLKMANN, Andreas; MÜLLER, Klaus; WAGNER, Sabine. Low-cost IoT irrigation system with MQTT communication and mobile interface. \textbf{Computers and Electronics in Agriculture}, v. 195, p. 106-118, 2022.

\bibitem{medeiros2021}
MEDEIROS, Paulo R.; NASCIMENTO, Carla M.; XAVIER, Roberto A. Sistema de irrigação automatizada para plantas domésticas baseado em Arduino e aplicação web. \textbf{Revista de Engenharia e Pesquisa Aplicada}, v. 6, n. 2, p. 45-56, 2021.

\bibitem{lima2019}
LIMA, Sandra J.; MOREIRA, Antônio C.; SILVA, Eduardo P. Automação da irrigação na agricultura familiar: revisão da literatura e perspectivas futuras. \textbf{Revista Brasileira de Agricultura Irrigada}, v. 13, n. 4, p. 3542-3558, 2019.

\bibitem{costa2020}
COSTA, Marina L.; FERREIRA, João P.; RODRIGUES, Ana C. Interfaces adaptativas em sistemas IoT: design centrado no usuário para automação residencial. \textbf{Interacting with Computers}, v. 32, n. 3, p. 298-312, 2020.

\bibitem{brasileiro2021}
BRASILEIRO, Ricardo M.; SOUZA, Patricia N. Calibração de sensores capacitivos para diferentes tipos de substrato: metodologia e validação experimental. \textbf{Measurement}, v. 178, p. 109-117, 2021.

\bibitem{martinez2022}
MARTINEZ, Elena S.; THOMPSON, James R.; GARCIA, Miguel A. User experience evaluation in home automation systems: metrics and methodologies. \textbf{International Journal of Human-Computer Studies}, v. 163, p. 102-115, 2022.

\end{thebibliography}

\end{document}
